\section{Related Work}

\subsection{Visual Document Retrieval}

Visual document retrieval has advanced rapidly by encoding document page images with vision-language models (VLMs) and matching queries via late interaction over token-level embeddings \citep{faysse2024colpali,khattab2020colbert,santhanam-etal-2022-colbertv2}. ColPali is a representative OCR-free framework in this line: it produces multi-vector representations from page screenshots and applies late-interaction scoring, avoiding complex text extraction and ingestion pipelines while achieving strong retrieval performance on visually rich documents \citep{faysse2024colpali}. Subsequent work has strengthened the visual backbone, for example through ColQwen \citep{faysse2024colpali}, explored more efficient retrieval architectures such as DSE \citep{ma2024dse}, and integrated visual retrievers into end-to-end visual RAG systems such as VisRAG \citep{yu2025visrag}. To reduce per-page indexing cost, efficient variants such as Light-ColPali \citep{ma2025lightcolpali} and DocPruner \citep{yan2025docpruner} further reduce the number of document-side vectors. However, these methods still operate at the \emph{page} level: they index and rank individual pages rather than complete documents. In contrast, this work targets \emph{document-level} visual retrieval and represents an entire multi-page document with a single composed encoding.

\subsection{Retrieval Granularity and Benchmarks}

Existing benchmarks for visual document retrieval are largely page-centric. ViDoRe \citep{faysse2024colpali,mace2025vidore2} evaluates cross-document retrieval where the retrieval unit is a single page. MMDocIR \citep{dong2025mmdocir} focuses on intra-document page-level and layout-level retrieval. VisDoMBench \citep{suri2025visdombench} studies multi-document question answering rather than document ranking. As a result, current benchmarks do not provide a standard setting for \emph{document-level} visual retrieval, where the system must rank whole documents from a corpus of multi-page documents.

A natural alternative is to reuse ideas from text retrieval by indexing smaller units and aggregating them to the document level. Passage-to-document aggregation has been widely studied in text retrieval, where methods such as MaxP aggregate passage scores into document scores \citep{dai2019deeper,li2023parade}. Multi-page document understanding has also been extensively explored for long-document question answering and reasoning, for example in Hi-VT5 \citep{tito2023hivt5}, DUDE \citep{Van_Landeghem_2023_ICCV}, M-LongDoc \citep{chia-etal-2025-longdoc}, and DocOwl~2 \citep{hu2025docowl2}. Nevertheless, these settings differ from document retrieval: the goal is typically to answer questions over a given document rather than to rank documents by relevance. To the best of our knowledge, there is still no benchmark or systematic study for cross-document \emph{document-level} visual retrieval with multi-positive relevance annotations.